% section: 解法の決定

\section{解法を選ぶ}\label{節：解法の決定}

  \chapterref{章：漸化式の解説}の解法を,式の特徴から選ぶ.
  「式全体」「\bookterm{sequence}{数列}の項」「付加項」のうち,目立つ特徴から変形を試す.
  複数の入口があってよく,変形後の式にも同じ見方を使える.

  \noindent\textbf{【解法を選ぶ手順】}\par
  次の4段階を繰り返す.チャートの出口は候補である.
  \begin{enumerate}
    \item \textbf{観察}：総和,分母,倍率,付加項など,目立つ特徴を見つける.
    \item \textbf{候補}：何を変えると,その特徴が消えるか,一定になるかを考える.
    \item \textbf{条件}：割る量の非零,\bookterm{logarithm}{対数}の真数の正値,初項による定数解の分岐を確かめる.
    \item \textbf{変形と評価}：新しい量の更新を一行計算する.簡単にならなければ,残った項から候補を選び直す.
  \end{enumerate}
  使える候補の中で,条件の見やすさと計算量を比べる.
  \bookterm{general}{一般項}が得られたら,\bookterm{initial}{初期条件}と元の更新式を確認する.
  \BookFigFlowchart

  \begingroup
  \clubpenalty=10000
  \widowpenalty=10000
  \medskip
  \noindent \textbf{【式全体を見る】}
  \begin{enumerate}
    \item \textbf{総和$S_n$}

      $S_{n+1}-S_n=a_{n+1}$を使い,添字をずらした式との差で総和を消す（総和型）.

    \item \textbf{項を含む分母}

      逆数を試す.\bookterm{constant}{定数係数}の1次式同士の商なら,定数を引いてから逆数を取る方法もある（分数型）.
      分母の非零を確かめる.

    \item \textbf{項数と添字の間隔}

      $a_{n+2}$と$a_n$の関係なら,奇数・偶数番目に分ける（非隣接2項間型）.
      また,$a_{n+2}=qa_n$は$a_{n+2}=0\cdot a_{n+1}+qa_n$とも書ける.
      定数係数で各項が1次の和なら,欠けた項を係数$0$で補い,隣接3項間・4項間の解法を試せる.

    \item \textbf{展開・因数分解・通分}

      整理すると基本型が現れることもある.約分で消す因子が零になる場合は別に調べる.

    \item \textbf{繰り返すまとまり}

      $a_n+n$と$a_{n+1}+(n+1)$が現れるなら,$b_n=a_n+n$を追う.

    \item \textbf{既知の公式}

      $a_n^2-2$なら,積が$1$の二数の和を試す（基本対称式型）.
      $2a_n^2-1$や$2a_n/(1-a_n^2)$なら,\bookterm{trigonometric}{三角関数}の倍角公式を試す.
      $\frac12(a_n+c/a_n)$の平均なら,\bookterm{newton}{ニュートン法}型を考える.
      初項と置換の条件も確かめる.

    \begin{samepage}
    \item \textbf{複数の\bookterm{sequence}{数列}}\\*[3pt]
      定数係数の\bookterm{linear}{線形}連立漸化式なら,和と差などで別々の更新へ分ける（連立型）.
    \par
    \end{samepage}

    \item \textbf{整数部分}

      記号の位置を見る（ガウス型）.項の値なら整数性や余り,添字なら偶奇や群,付加項なら同じ値を取る範囲に注目する.
      \bookterm{floor-function}{床関数}を外すには,整数性と上下の不等式を確かめる.
  \end{enumerate}

  \medskip
  \noindent \textbf{【\bookterm{sequence}{数列}の項を見る】}

  係数については,主に$a_{n+1}=p_na_n+f(n)$の形に整理できる場合を考える.
  \begin{enumerate}
    \item \textbf{係数の変化}

      \begin{itemize}[beginpenalty=10000]
        \item 一定なら等比型や定数項を消す方法を試す.$1$なら\bookterm{difference}{階差}を見る.
        \item $n$に依存するなら積を調べる（\bookterm{ratio}{階比}型）.
              付加項もあれば,数列で割って倍率をそろえる.例えば係数$(n+1)/n$なら$b_n=a_n/n$が候補になる.
      \end{itemize}
      付加項も併せて見る.

    \item \textbf{累乗の指数}

      \begin{itemize}[beginpenalty=10000]
        \item 2次以上なら因数分解や$a_n^2$などの置換を試す.
        \item 指数が定数でも$n$に依存しても,\bookterm{logarithm}{対数}で係数へ取り出せる場合がある（対数型）.
      \end{itemize}
      実数の対数には真数の正値が必要である.

    \item \textbf{項どうしの積と商}

      因数分解や置換,対数を試す.対数は積・商・累乗を和と差や定数倍へ変えるが,和は分解できない.
  \end{enumerate}

  \medskip
  \noindent \textbf{【付加項を見る】}

  $a_{n+1}=p_na_n+f(n)$で,未知の項を含まずに足される$f(n)$を「付加項」と呼ぶ.
  \begin{enumerate}
    \item \textbf{付加項の有無と定数}

      なければ係数の積を取る（等比型・階比型）.
      係数と付加項が一定なら定数を引く.ただし$a_{n+1}=a_n+d$はそのまま等差型で解ける.

    \item \textbf{添字への依存}

      \begin{itemize}[beginpenalty=10000]
        \item \textbf{指数}：$cq^n$なら,$q\ne0$を確かめ,$q^n$などで割る.
        \item \textbf{多項式}：係数が定数なら多項式を引く（階差等比複合型）.
        \item \textbf{多項式と指数の積}：$n2^n$なら指数部分で割り,残った多項式を扱う.
      \end{itemize}
      $n+2^n$のような和は,各部分を見る.

    \item \textbf{和を求める工夫}

      階差型へ帰着したら,$n(n+1)$などの連続整数の積や,隣り合う項の差へ直す.差なら途中の項が消える.
  \end{enumerate}

  \medskip
  \noindent \textbf{【補足：初項の手掛かり】}

  初項$1/(\sqrt6-\sqrt2)$を有理化すると$(\sqrt6+\sqrt2)/4=\cos15^\circ=\sin75^\circ$となる.
  三角関数を試す手掛かりだが,更新式の形と合わせて判断する.
  \par
  \endgroup

% 解法の選択を、計算の前に確認する短い練習。
\subsection{操作を選ぶ練習}\label{節：操作を選ぶ練習}
  \bookterm{general}{一般項}の計算の前に,新しい量の更新を確かめる.
  引く・割る・\bookterm{logarithm}{対数}を取る操作を比較する.

  \noindent\bookblocklink{operation-choice}{problem}{answer}{【問題】}\par\nobreak
  \begin{enumerate}[format=\bookitemlink{operation-choice}{problem}{answer}]
    \item $a_1=1$,\,$a_{n+1}=3a_n+2n+1$とする.
          定数を引くだけで付加項を消せるか.
          $g(n)=un+v$を引く場合の$u,v$を求め,\,$b_n=a_n-g(n)$の更新式と初項を示せ.
    \item $a_1=2$,\,$a_{n+1}=\frac{n+2}{n}a_n+(n+1)(n+2)$とする.
          $b_n=\frac{a_n}{n}$と$c_n=\frac{a_n}{n(n+1)}$を候補とし,それぞれの更新式を求めよ.
          どちらで$a_n$の係数を$1$にそろえられるか.
    \item $a_1=-2$,\,$a_{n+1}=2^{2n+1}a_n$とする.
          $b_n=\log_2a_n$は使えるか.
          \bookterm{logarithm}{対数}を使わず,\,$c_n=\frac{a_n}{2^{n^2}}$の更新式と初項を求めよ.
  \end{enumerate}

  \noindent\bookblocklink{operation-choice}{answer}{problem}{【解答と選択の理由】}\par\nobreak
  \begin{enumerate}[format=\bookitemlink{operation-choice}{answer}{problem}]
    \item $a_n-\alpha$では付加項の$2n$が残る.
          $g(n+1)=3g(n)+2n+1$から
          \[
            u=3u+2,\qquad u+v=3v+1
          \]
          となり,\,$u=v=-1$である.
          従って$b_n=a_n+n+1$により$b_{n+1}=3b_n$,\,$b_1=3$を得る.
          引く量を定数から一次式へ変えたことで,付加項を消せた.

    \item 候補ごとに一歩だけ更新すると,
          \[
            b_{n+1}=\frac{n+2}{n+1}b_n+n+2
          \]
          と
          \[
            c_{n+1}=c_n+1,\qquad c_1=1
          \]
          を得る.
          $h_n=n(n+1)$なら
          \[
            \frac{h_{n+1}}{h_n}=\frac{n+2}{n}
          \]
          が元の倍率と一致する.
          だからこの因子で割る.
          $n\ge1$なので分母は非零であり,\,$c_n=n$から$a_n=n^2(n+1)$も得られる.

    \item 正の倍率で更新するので,各項は負であり,\,$\log_2a_n$は実数として定義できない.
          一方,\,$2n+1=(n+1)^2-n^2$から
          \[
            c_{n+1}=c_n,\qquad c_1=-1
          \]
          となる.従って$a_n=-2^{n^2}$を得る.
          \bookterm{normalize}{正規化}で割る$2^{n^2}$は常に非零であり,元の項が正である必要はない.
          初項が正の比較例は\bookproblemref{practice-basic}{42}{標準問42}へ.
  \end{enumerate}

  \noindent\textbf{【同じ式を二つの方法で見る】}\par
  第1問では,一次式$-n-1$を引くと等比型になり,$a_n=3^n-n-1$を得る.
  一方,$c_n=\dfrac{a_n}{3^{n-1}}$と置いてもよい.この場合は
  \[
    c_{n+1}-c_n=\frac{2n+1}{3^n}
  \]
  となり,残った和を計算して元へ戻す.
  どちらも$n\ge1$で使える.ここでは一次式がすぐに決まり,引けば和の計算も要らない.
  引く量が見つからなければ,倍率から割る因子を作れる.残る計算を比べて選ぶ.

  \noindent\textbf{【倍率をそろえる仕組み】}\par
  第2問は\sectionref*{節：階比型}の操作を付加項がある場合へ広げたものだ.
  一般に
  \[
    a_{n+1}=p_na_n+q_n
  \]
  に対して,既知の非零の\bookterm{sequence}{数列}$h_n$が$h_{n+1}=p_nh_n$を満たすなら,
  \[
    b_n=\frac{a_n}{h_n},\qquad
    b_{n+1}-b_n=\frac{q_n}{h_{n+1}}
  \]
  となる.倍率をそろえた後に,\bookterm{difference}{階差}を足し戻す.
  $p_n\ne0$なら,\,$h_1=1$,\,$h_n=\prod_{k=1}^{n-1}p_k$を候補にできる（空積は$1$）.
  零の係数はその段階を別に調べる.残った和が簡単に計算できるとは限らない.

